\chapter{Resultados} \label{resultados}


% \section{Uma associação}
%
% % Gera 3 parágrafos com texto aleatório, apenas para exemplo. Apague o comando para remover tal texto.
% \lipsum[1-3]
%
%
% \subsection{Motivos para abortar uma associação}
%
% %lista não numerada
% \begin{itemize}
% 	\item primeiro item;
% 	\item segundo item;
% 	\item terceiro item;
% 	\item quarto item.
% \end{itemize}
%
% \begin{figure}[h]
% 	\centering
% 	\includegraphics[scale=0.8]{images/red.png}
% 	\caption{Exemplo de imagem}
% 	\label{imagem1}
% \end{figure}
%
%
% Se durante o processo de configuração de uma associação for recebido como payload um hostname e esse hostname não puder ser resolvido em um tempo hábil deve se enviar um abort com a causa de erro de endereço não resolvido. Veja exemplo na Figura \ref{imagem1}.
%
%
% % Gera 4 parágrafos com texto aleatório, apenas para exemplo. Apague o comando para remover tal texto.
% \lipsum[1-4]
%
% A Tabela \ref{tabela1} a seguir apresenta os tipos de \textit{chunk} de um pacote do protocolo.
%
% \begin{table}[ht!]
%   \begin{center}
%   \setlength{\belowcaptionskip}{10pt} % espao entre caption e tabela
%   \footnotesize {
%       \begin{tabular}{|p{4cm}|p{9cm}|}
% 	  \hline
% 	  \textbf{Nome} & \textbf{Função} \\
% 	  \hline
% 	  Iniciar & Usado para iniciar uma associação \\
% 	  \hline
% 	  Confirmacao & Segunda mensagem de uma configuração de uma associação\\
% 	  \hline
% 	  Mensagem & Terceira mensagem de uma configuração de uma associação\\
% 	  \hline
% 	  Cookie & Quarta mensagem de uma configuração de uma associação\\
% 	  \hline
% 	  Dados & Dados da aplicação\\
% 	  \hline
%       \end{tabular}
%   }
%   \caption{Tipos de \textit{chunk} de um pacote SCTP}
%   \label{tabela1}
%   \end{center}
% \end{table}
